\subsection{Boxplots, Scatterplots, and Relationship Summaries}
\label{subsec:boxplots-scatterplots-and-relationship-summaries}
% Source integration: attached screenshots from Chapter 10, "Summarizing Data."

Boxplots and scatterplots are compact exploratory displays.  A boxplot summarizes one batch or compares several batches; a scatterplot displays relationships between paired variables.

\begin{method} \label{meth:boxplot-construction} \hlt{Boxplot Construction}\\
A standard boxplot is built from the median, lower quartile $Q_1$, upper quartile $Q_3$, and $\operatorname{IQR}=Q_3-Q_1$:
\begin{enumerate}[label=(\roman*)]
\setlength{\itemsep}{0pt}
\item Draw a box from $Q_1$ to $Q_3$ and mark the median inside the box.
\item Draw whiskers to the most extreme observations lying within $1.5\operatorname{IQR}$ of the nearer quartile.
\item Mark observations beyond the whiskers individually as possible outliers.
\end{enumerate}
\end{method}

\begin{figure}[H]
  \centering
  % Source: supplied Chapter 10 screenshots on boxplot construction.
% Schematic redraw, not a reproduction of the scanned figures.
\begin{tikzpicture}[
  x=0.9cm,
  y=0.40cm,
  >=Stealth,
  line cap=round,
  line join=round,
  every node/.style={font=\scriptsize},
  axis/.style={->, gray!70, line width=0.45pt},
  box/.style={draw=blue!70!black, fill=blue!8, line width=0.7pt},
  medianmark/.style={blue!70!black, line width=0.8pt},
  whisker/.style={blue!70!black, line width=0.7pt},
  outlier/.style={black, fill=black, inner sep=0pt, minimum size=2pt, circle},
  guide/.style={gray!35, dashed, line width=0.35pt},
  label/.style={text=black!78}
]
  \draw[axis] (0,0) -- (0,8.9) node[black!78,above] {value};

  \draw[guide] (.85,3) -- (3.35,3);
  \draw[guide] (.85,4.05) -- (3.35,4.05);
  \draw[guide] (.85,5.8) -- (3.35,5.8);

  \draw[box] (1.35,3) rectangle (2.55,5.8);
  \draw[medianmark] (1.35,4.05) -- (2.55,4.05);
  \draw[whisker] (1.95,1.15) -- (1.95,3);
  \draw[whisker] (1.6,1.15) -- (2.3,1.15);
  \draw[whisker] (1.95,5.8) -- (1.95,6.85);
  \draw[whisker] (1.6,6.85) -- (2.3,6.85);
  \foreach \y in {7.55,8.15,8.65}
    \node[outlier] at (1.95,\y) {};

  \node[label, right] at (3.42,5.8) {$Q_3$};
  \node[label, right] at (3.42,4.05) {median};
  \node[label, right] at (3.42,3) {$Q_1$};
\end{tikzpicture}

  \caption{A boxplot displays center, central spread, skewness, and possible outliers.}
  \label{fig:eda-boxplot-schematic}
\end{figure}

\begin{remark} \label{rem:boxplots-compare-batches} \hlt{Comparing Batches with Boxplots}\\
Parallel boxplots are useful when many batches must be compared at once.  Differences in medians suggest shifts in location; differences in box lengths suggest changes in central spread; asymmetric whiskers and one-sided outliers suggest skewness or tail imbalance.
\end{remark}

\begin{definition} \label{def:scatterplot-and-correlation-summaries} \hlt{Scatterplot, Pearson Correlation, and Rank Correlation}\\
For paired data $(x_1,y_1),\ldots,(x_n,y_n)$:
\begin{enumerate}[label=(\roman*)]
\setlength{\itemsep}{0pt}
\item A \hlt{scatterplot} places the points in the $xy$-plane to display association, curvature, clusters, and outliers.
\item The \hlt{Pearson correlation coefficient} is
\[
r=\frac{\sum_{i=1}^n(x_i-\overline{x})(y_i-\overline{y})}
{\sqrt{\sum_{i=1}^n(x_i-\overline{x})^2}\sqrt{\sum_{i=1}^n(y_i-\overline{y})^2}}.
\]
\item The \hlt{rank correlation coefficient} is the Pearson correlation of the ranks of the $x_i$ and the ranks of the $y_i$.
\end{enumerate}
Pearson correlation measures linear association on the displayed scales.  Rank correlation is invariant under monotone transformations and is less sensitive to extreme numerical values.
\end{definition}

\begin{method} \label{meth:scatterplot-diagnostic-workflow} \hlt{Scatterplot Diagnostics}\\
When using scatterplots:
\begin{enumerate}[label=(\roman*)]
\setlength{\itemsep}{0pt}
\item Plot the raw variables first and look for outliers, curvature, clusters, and changing spread.
\item Try scientifically meaningful transformations, such as logarithms for variables spanning several orders of magnitude.
\item Use correlation only after checking whether a linear summary is visually appropriate.
\item For more than two variables, inspect a scatterplot matrix before relying on pairwise numerical summaries alone.
\end{enumerate}
\end{method}

\begin{example} \label{ex:scale-can-hide-relationship-in-scatterplot} \hlt{Changing Scale Can Reveal Structure}\\
If one coordinate varies over several orders of magnitude, a raw scatterplot can compress most points near an axis and hide a relationship.  Plotting against the logarithm of that coordinate can reveal an approximately monotone or linear pattern on the transformed scale.  The correlation before and after such a transformation need not agree, because Pearson correlation measures linear association on the chosen coordinates.
\end{example}

\begin{figure}[H]
  \centering
  % Source: supplied Chapter 10 screenshots on scatterplots and logarithmic scale changes.
% Schematic redraw, not a reproduction of the scanned figures.
\begin{tikzpicture}[
  x=0.72cm,
  y=0.36cm,
  >=Stealth,
  line cap=round,
  line join=round,
  every node/.style={font=\scriptsize},
  axis/.style={->, gray!70, line width=0.45pt},
  point/.style={black, fill=black, inner sep=0pt, minimum size=1.9pt, circle},
  trend/.style={blue!70!black, line width=0.95pt}
]
  \begin{scope}
    \draw[axis] (0,0) -- (5.2,0) node[black!78,right] {$x$};
    \draw[axis] (0,0) -- (0,9.2) node[black!78,above] {$y$};
    \foreach \x/\y in {
      .2/7.9,.32/6.1,.42/8.5,.52/5.2,.63/7.0,.75/4.3,
      .9/6.2,1.1/4.9,1.35/5.4,1.65/4.0,4.35/2.75
    } \node[point] at (\x,\y) {};
  \end{scope}

  \begin{scope}[xshift=5.1cm]
    \draw[axis] (0,0) -- (5.2,0) node[black!78,right] {$\log x$};
    \draw[axis] (0,0) -- (0,9.2) node[black!78,above] {$y$};
    \foreach \x/\y in {
      .35/8.45,.75/7.9,1.15/7.35,1.55/6.65,1.95/5.95,
      2.35/5.35,2.75/4.75,3.15/4.2,3.55/3.65,3.95/3.05
    } \node[point] at (\x,\y) {};
    \draw[trend] (.25,8.45) -- (4.25,2.75);
  \end{scope}
\end{tikzpicture}

  \caption{A change of scale can turn a compressed scatterplot into a readable relationship display.}
  \label{fig:eda-scatterplot-scale-change}
\end{figure}

\begin{remark} \label{rem:exploratory-displays-before-inference} \hlt{Exploration Before Formal Inference}\\
Exploratory displays do not replace confidence intervals, tests, or models.  They often determine which formal question is sensible: whether a model should include a transformation, whether an apparent relationship is driven by a few points, whether batches have comparable spread, and whether a data-recording problem is plausible.
\end{remark}
